\subsection{Experimental Details: Language Transfer}
\label{app:language-transfer-details}

In \Cref{sec:lang-transfer} we measure how much training on \emph{source language} $s$ is beneficial for the test loss on the \emph{target language} $t$.
A language transfer score tells a practitioner the extent to which training with each language would help or hinder their target language $t$.
There are two ways in which we can measure this language transfer:

\begin{enumerate}[noitemsep]

    \item \textsc{Bilingual Transfer Score.} This score measures how many training tokens a 50/50 ($s$, $t$) bilingual model needs, relative to a monolingual target language $t$ baseline, to reach the same validation loss $L_t$. A positive score indicates co-training with the source language has positive transfer, whereas a negative score indicates it hinders convergence.

    \item \textsc{Finetuning Adaptation Score.} This score captures the effect on the target language's loss $L_t$, when a massively multilingual model is finetuned only on the source language $s$. 

\end{enumerate}

\subsubsection{\bts{} (\btsabb)}

The \bts{} asks how data-efficient a bilingual learner is, relative to a purely monolingual one, on some target language $t$.
To estimate this, we pretrain two models, a monolingual one on $t$ and a bilingual one that samples tokens equally (50/50) from the source and target languages ($s$, $t$).
We then measure how many more training tokens it takes the bilingual model to attain the same validation loss $L_t$ as the monolingual model.
This ``distance'' between the learning curves is measured throughout training, at different reference horizons---or, number of tokens trained on.
In \Cref{fig:language-transfer-heatmap} we use models with 2B parameters, and a reference horizon of $d\!=\!42B$ tokens, as it roughly equates to $10,000$ pretraining steps for our 2B parameter models, and learning curves have stabilized.
In \Cref{fig:language-transfer-nd} we measure how the \btsabb{} varies by model size, and by number of training tokens seen.

We are able to compute the bilingual transfer scores, across every combination of pairings between 10 languages: English, Russian, Spanish, French, German, Portuguese, Chinese, Vietnamese, Japanese, and Hindi.
These languages were chosen to give a distribution of high- and mid-resource languages from a variety of language families.
This totals 10 monolingual experiments (one for each language), and $C(10,2)=45$ bilingual experiments, which provide the results for a 10x10 grid of Bilingual Transfer Scores. 

In more detail, let $d_{mono}$ be the number of tokens trained on language $t$ by a monolingual model. This model attains test loss of $L_{t}(d_{mono})$ after $d_{mono}$ tokens on language $t$.  
We record the number of training tokens $d_{bi}$ at which the bilingual model reaches the same loss, $L_{t}(d_{\text{bi}}) = L_{t}(d_{mono})$.
We define the \textbf{Bilingual Transfer Score} (BTS) as the signed step surplus

\[
\mathrm{BTS}_{s\!\rightarrow t}  = - \frac{\sigma_{\text{bi}}(L_{\text{t}}(d_{\mathrm{mono}})) - 2d_{\mathrm{mono}}}{d_{\mathrm{mono}}}
\]

where $d_{\text{bi}}=\sigma_{\text{bi}}(L_{\text{t}}(d_{\mathrm{mono}}))$. $\sigma_{\text{bi}}$ maps a loss value to the number of steps taken by the bilingual model to reach that loss.
In short, this score measures the relative training efficiency of the bilingual model compared to a monolingual model at reaching the same loss level for language $t$. 
Subtracting $2d_{\mathrm{mono}}$ and dividing by $d_{\mathrm{mono}}$ simply centers the value on 0, and forces positive transfer to give a positive value.
A value of 0, implies \emph{neutral} transfer, as the multilingual model requires exactly the same number of steps in the target language $t$ as the monolingual model, a value below 0 implies \emph{negative} transfer, as the multilingual model is slower to reach the loss $L_{\text{t}}(d_{\mathrm{mono}})$, and a value above 0 implies \emph{positive} transfer, as the bilingual model reaches the target loss in $<2d_{\mathrm{mono}}$.
Note that for \Cref{fig:language-transfer-heatmap} these scores are anchored on 2B parameter models, trained for $d\!=\!42B$ tokens, for consistency.
But as model size increases, or tokens trained decreases, the language transfer scores improve monotonically.

\subsubsection{\fas{} (\fasabb)}

Complementary to \btsabb{}, the \fas{} measures the \emph{instantaneous impact} of continued exposure to a single language on all others while holding model capacity fixed.  
Specifically, it measures how finetuning a massively multilingual model on source language $s$ affects the performance of target language $t$.
Because we only need to train on each source language once, and evaluate on all languages, it is less computationally expensive than \btsabb{}, which requires training on every pair of languages.
As such, we are able to derive $38*38=1444$ ($s$, $t$) comparisons from 38 finetuning experiments.

We start from a shared multilingual checkpoint: a \unimax{} pretrained model, trained for $1T$ tokens.
For each source language $s$, we continue pre-training (for simplicity we denote this as \emph{finetune}) exclusively on this language, and evaluate at regular intervals on every target language's $t\!\in\!\mathcal{L}$ validation set.  
Let $L_{s\!\rightarrow\!t}(d)$ be the validation loss on $t$ after $d$ additional updates on $s$, and let $L^{unimax}_t$ denote the baseline loss of the shared \unimax{} checkpoint, before finetuning has begun.  
The \fas{} (\fasabb) aggregates the reduction in loss over a common time window $[0,d_{\max}]$:
\[
\mathrm{FAS}_{s\!\rightarrow t}
=\frac{1}{d_{\max}}
\int_{0}^{d_{\max}}
\!\bigl[L^{unimax}_t - L_{s\!\rightarrow\!t}(d)\bigr]\,
\mathrm{d}d.
\]
Positive values indicate that exposing the model to $s$ accelerates learning or yields immediate gains on $t$ (``helpful transfer''), whereas negative values capture negative interference. Normalising by $d_{\max}$ makes the score comparable across language pairs and training horizons.

Because finetuning on language $s$ can cause unpredictable behavior in the validation loss of a target language $t$ (it might immediately increase, or decrease then eventually increase), we start with deriving a series of features about the finetuning loss curve. 
We calculate the following features for a finetuning loss curve:

\begin{itemize}%[noitemsep]

    \item \textbf{Baseline Loss Deviation ($\delta_{s\rightarrow t}$).} 
    For each language pair $s\neq t$ we quantify the deviation in validation loss between the baseline $L^{unimax}_t$ and the loss $L_{s \rightarrow t}$ on the final step: $$\delta_{s\rightarrow t} = L_{s \rightarrow t}(d_\text{max}) - L^{unimax}_t(d_\text{max})$$
    $\delta_{s\rightarrow t}<0$ means fine–tuning on $s$ helps\/ $t$;
    $\delta_{s\rightarrow t}>0$ means it hurts.
    
    \item \textbf{Adaptation Gain ($g_{l}$).} This measures the area normalized reduction in loss on language $l$ from finetuning on language $l$. This allows us to capture the relative learning dynamics of both the source and target language.
    $$g_{l} = \frac{1}{d_{max}} \int_{0}^{d_{\text{max}}} \bigl[\,L^{unimax}_l - L_{s \rightarrow t}(d)\,\bigr]\; \mathrm{d}s.$$
    $g_l>0$ indicates net improvement; large values imply that $l$
    benefits greatly from additional supervision.

\end{itemize}

For each language transfer pair ($s$, $t$), we compose the following feature vector $x_{s \rightarrow t}$. 
As languages differ in intrinsic difficulty, we normalize each feature. 
%
\begin{align}
\tilde{g}_l &=
\frac{g_l - \mu_g}{\sigma_g},
&
\tilde{\delta}_{s\rightarrow t} &=
\frac{\delta_{s\rightarrow t}-\mu_{\delta,t}}{\sigma_{\delta,t}},
\label{eq:zscore}
\end{align}
%
For adaptation gain, we use the global mean and standard deviation $(\mu_g,\sigma_g)$, and for baseline loss deviation we compute the mean and standard deviation $(\mu_{\delta,t},\sigma_{\delta,t})$ across all source languages $s$ for each fixed target language $t$.

\subsection{Estimating \bts{} (\btsabb)}
\label{app:lang-transfer-bts-estimation}

% From there we compute the feature vector $x_{s \rightarrow t}$ using the 90 \bts experiments.

In the \bts{} experiments we empirically measured language transfer scores
${y}_{s\rightarrow t}^{\text{true}}$ for a subset
$\mathcal{P}_{\text{obs}}\subset\mathcal{L}\!\times\!\mathcal{L}$. 
We construct a training set from these 90 experiments, using them as the gold labels.
%
\begin{equation}
\mathbf{x}_{s\rightarrow t}
=
\bigl[
\tilde{g}_s,\;
\tilde{g}_t,\;
\tilde{g}_s-\tilde{g}_t,\;
\tilde{\delta}_{s\rightarrow t}
\bigr]^{\!\top},
\qquad
y_{s\rightarrow t}=y^{\text{true}}_{s\rightarrow t},
\label{eq:features}
\end{equation}
%
% excluding trivial self–pairs $(s=s)$.  
% Missing $\tilde{\delta}_{s\rightarrow t}$ values (rare) are imputed with the column mean.

We fit a random-forest regressor $\varphi_{\boldsymbol{\theta}}:\mathbb{R}^4\!\rightarrow\!\mathbb{R}$ with $300$ trees and unlimited depth.  
The predictive performance is assessed by $k$-fold cross-validation
($k\!=\!5$), obtaining an $R^{2}=0.85$ and Spearman correlation of $\rho=0.88$.
After training on all observed pairs, $\varphi$ is invoked on every unlabelled $(s,t)\!\in\!\mathcal{L}\!\times\!\mathcal{L}$, yielding predicted \btsabb{} values, that estimate the language transfer shown in \Cref{fig:language-transfer-heatmap}.



